\section{Yorug‘lik kvantlari. Fotonlar}



\begin{enumerate}

\item Jismlarning nurlanishi to‘g‘risidagi Plank gipotezasi qanday?
\begin{enumerate}[label=\Alph*)]
\item Jismlarning nurlanishi temperaturaga bog‘liq ravishda ortib boradi.
\item Jismlarning nurlanishi uzluksiz ravishda davom etib turadi.
\item Jismlarning nurlanishi doimiy jarayondir.
\item Jismlarning nurlanishi uzluksiz emas, balki alohida ulushlar (kvantlar) tarzida chiqariladi.
\end{enumerate}

\item Plank doimiysining birligini ko‘rsating.
\begin{enumerate}[label=\Alph*)]
\item J$\cdot$m/s.
\item J.
\item J/s.
\item J$\cdot$s.
\end{enumerate}

\item Foton energiyasini aniqlaydigan ifodani ko‘rsating:
1) $E = h\nu$; 2) $E = hc/\lambda$; 3) $E = h\nu/c$; 4) $E = h\lambda/c$;
\begin{enumerate}[label=\Alph*)]
\item 2 va 3.
\item 1 va 4.
\item 1.
\item 1 va 2.
\end{enumerate}

\item Chastotasi $\nu$ bo‘lgan fotonning energiyasiga mos ifodani toping.
\begin{enumerate}[label=\Alph*)]
\item $h\nu$.
\item $h\nu c$.
\item $h\nu c^2$.
\item $h\nu/c$.
\end{enumerate}

\item Energiyasi E ga teng bo‘lgan fotonning chastotasi nimaga teng?
\begin{enumerate}[label=\Alph*)]
\item $Eh$.
\item $E/h$.
\item $E/c$.
\item $Eh/c^2$.
\end{enumerate}

\item Energiyasi $6,4\cdot 10^{-18}$ J bo‘lgan foton energiyasini eV larda ifodalang.
\begin{enumerate}[label=\Alph*)]
\item 0,4.
\item 4.
\item 40.
\item 400.
\end{enumerate}

\item $1,76\cdot 10^{-14}$ J necha keV?
\begin{enumerate}[label=\Alph*)]
\item 11.
\item 1,1.
\item 110.
\item 0,11.
\end{enumerate}

\item $19,2\cdot 10^{-10}$ J necha MeV?
\begin{enumerate}[label=\Alph*)]
\item 12.
\item 1,2.
\item 120.
\item 12000.
\end{enumerate}

\item 30 MeV necha Joul?
\begin{enumerate}[label=\Alph*)]
\item $9,2\cdot 10^{-10}$.
\item $19,2\cdot 10^{-10}$.
\item $1,2\cdot 10^{-11}$.
\item $4,8\cdot 10^{-12}$.
\end{enumerate}

\item Atom chastotasi $3,2\cdot 10^{16}$ Hz bo‘lgan nurlanish kvantini yutganda, atomning energiyasi necha eV ga o‘zgaradi? $h = 6,6\cdot 10^{-34}$ J$\cdot$s; 1 eV = $1,6\cdot 10^{-19}$ J.
\begin{enumerate}[label=\Alph*)]
\item 160.
\item 150.
\item 264.
\item 132.
\end{enumerate}

\item Kislorod atomini ionlashtirish uchun 14 eV energiya zarur. Ionlashtiruvchi nurlanishning chastotasini toping.
\begin{enumerate}[label=\Alph*)]
\item $34\cdot 10^{15}$ Hz.
\item $3,4\cdot 10^{15}$ Hz.
\item $2\cdot 10^{16}$ Hz.
\item $17\cdot 10^{15}$ Hz.
\end{enumerate}

\item Kislorod atomini ionlashtirish uchun 16,5 eV energiya kerak. Ionlashtiruvchi nurlanishning chastotasi necha Gers? $h = 6,6\cdot 10^{-34}$ J$\cdot$s, 1 eV = $1,6\cdot 10^{-19}$ J.
\begin{enumerate}[label=\Alph*)]
\item $3\cdot 10^{15}$.
\item $5\cdot 10^{15}$.
\item $4\cdot 10^{15}$.
\item $4\cdot 10^{16}$.
\end{enumerate}

\item Agar chastotasi $3,4\cdot 10^{15}$ Hz bo‘lgan nurlar kislorod atomlarini ionlashtirishi mumkin bo‘lgan eng kichik chastotali nurlar bo‘lsa, atomlarni ionlashtirish energiyasi qanday? $h = 6,6\cdot 10^{-34}$ J$\cdot$s.
\begin{enumerate}[label=\Alph*)]
\item 0,14 eV.
\item 140 eV.
\item 1,4 keV.
\item 14 eV.
\end{enumerate}

\item Foton massasi formulasini ko‘rsating.
\begin{enumerate}[label=\Alph*)]
\item $Ec^2$.
\item $h\nu$.
\item $h\nu/c$.
\item $h\nu/c^2$.
\end{enumerate}

\item Fotonning impulsi $3,31\cdot 10^{-27}$ kg$\cdot$m/s bo‘lsa, unga mos keladigan to‘lqin uzunligi qanday (nm)? $h = 6,62\cdot 10^{-34}$ J$\cdot$s.
\begin{enumerate}[label=\Alph*)]
\item 200.
\item 331.
\item 400.
\item 500.
\end{enumerate}

\item Chastotasi $3\cdot 10^{15}$ Hz bo‘lgan fotonning impulsini aniqlang (kg$\cdot$m/s). $h = 6,63\cdot 10^{-34}$ J$\cdot$s.
\begin{enumerate}[label=\Alph*)]
\item $2,21\cdot 10^{-19}$.
\item $2,21\cdot 10^{-25}$.
\item $2\cdot 10^{-19}$.
\item $6,63\cdot 10^{-27}$.
\end{enumerate}

\item Absolut sindirish ko‘rsatkichi n bo‘lgan shaffof muhitda fotonning impulsi nimaga teng? ($\nu$, $\lambda$ -- fotonning shu muhitdagi chastotasi va to‘lqin uzunligi).
\begin{enumerate}[label=\Alph*)]
\item $nh\nu/c$.
\item $nh\nu$.
\item $h\lambda/n$.
\item $h\nu/nc$.
\end{enumerate}

\item Oq yorug‘lik bilan qizil, sariq, yashil, ko‘k rangdagi shishalarning qaysi biri yoritilganda, o‘tgan nur fotonining impulsi eng katta bo‘ladi?
\begin{enumerate}[label=\Alph*)]
\item yashil.
\item ko‘k.
\item qizil.
\item sariq.
\end{enumerate}

\item Impulsi $v = 10^4$ m/s tezlikda harakatlanayotgan elektronning impulsiga teng bo‘lgan fotonning to‘lqin uzunligi qanday (m)? $h = 6,62\cdot 10^{-34}$ J$\cdot$s, $m = 9,1\cdot 10^{-31}$ kg.
\begin{enumerate}[label=\Alph*)]
\item $7,3\cdot 10^{-8}$.
\item $7,3\cdot 10^{-6}$.
\item $5,5\cdot 10^{-8}$.
\item $7,3\cdot 10^{-5}$.
\end{enumerate}

\item Birinchi fotonning energiyasi ikkinchi fotonnikidan 2 marta katta bo‘lsa, birinchi foton impulsi ikkinchi foton impulsidan necha marta katta bo‘ladi?
\begin{enumerate}[label=\Alph*)]
\item 1/4.
\item 1/8.
\item 2.
\item 4.
\end{enumerate}

\item $6,62\cdot 10^{-7}$ m to‘lqin uzunlikda ishlayotgan lazer yorug‘lik dastasining quvvati 0,1 W. Lazer 3 s da nechta foton nurlaydi? $h = 6,62\cdot 10^{-34}$ J$\cdot$s.
\begin{enumerate}[label=\Alph*)]
\item $5\cdot 10^{17}$.
\item $1\cdot 10^{18}$.
\item $2,5\cdot 10^{21}$.
\item $5\cdot 10^{21}$.
\end{enumerate}

\item Cho‘g‘lanma lampa nurlarining o‘rtacha to‘lqin uzunligi $1,2\ \mu$m bo‘lsa, quvvati 200 W bo‘lgan lampa bir sekundda nechta foton nurlaydi? Iste’mol qilinadigan barcha quvvat nurlanishga sarflanadi va $h = 6,6\cdot 10^{-34}$ J$\cdot$s deb hisoblansin.
\begin{enumerate}[label=\Alph*)]
\item $4\cdot 10^{23}$.
\item $2,5\cdot 10^{21}$.
\item $1,8\cdot 10^{20}$.
\item $1,2\cdot 10^{21}$.
\end{enumerate}

\item Foydali quvvati 300 W bo‘lgan yorug‘lik manbai 662 nm uzunlikdagi yorug‘lik to‘lqini chiqaradi. Manbadan har sekundda nurlanayotgan fotonlarning soni qanday? $c = 3\cdot 10^8$ m/s; $h = 6,62\cdot 10^{-34}$ J$\cdot$s.
\begin{enumerate}[label=\Alph*)]
\item $5\cdot 10^{20}$.
\item $1\cdot 10^{20}$.
\item $1\cdot 10^{19}$.
\item $1\cdot 10^{21}$.
\end{enumerate}

\item Yorug‘lik manbai har sekundda $10^{20}$ ta foton chiqaradi. Nurlanishning to‘lqin uzunligi 662 nm bo‘lsa, manbaning quvvati qanday (W)? $h = 6,62\cdot 10^{-34}$ J$\cdot$s.
\begin{enumerate}[label=\Alph*)]
\item 10.
\item 20.
\item 30.
\item 60.
\end{enumerate}

\item Quvvati 50 W bo‘lgan yorug‘lik manbai har 2 s da o‘rtacha $2,5\cdot 10^{20}$ ta foton chiqaradi. Nurlanishning o‘rtacha to‘lqin uzunligi qanday ($\mu$m). $h = 6,63\cdot 10^{-34}$ J$\cdot$s.
\begin{enumerate}[label=\Alph*)]
\item 0,25.
\item 0,5.
\item 0,6.
\item 1.
\end{enumerate}

\item Yerga tushayotgan Quyosh nurlarining intensivligi $1,4\cdot 10^3$ W/m$^2$. Quyosh nurlarini batamom yutuvchi sirtga berilayotgan bosimni aniqlang (Pa).
\begin{enumerate}[label=\Alph*)]
\item $4,7\cdot 10^6$.
\item $1,4\cdot 10^3$.
\item $3\cdot 10^3$.
\item $4,7\cdot 10^{-6}$.
\end{enumerate}

\item $100\text{ cm}^2$ sirtga har minutda 63 J yorug‘lik energiyasi tushadi. Sirt tushuvchi nurlarni batamom qaytaradigan hol uchun yorug‘lik bosimi qanday (N/m$^2$) bo‘ladi?
\begin{enumerate}[label=\Alph*)]
\item $3,5\cdot 10^{-7}$.
\item $7\cdot 10^{-7}$.
\item $7\cdot 10^{-5}$.
\item $3,5\cdot 10^{-5}$.
\end{enumerate}

\item Qaytarish koeffitsenti 0,4 bo‘lgan sirtga intensivligi 1200 W/m$^2$ bo‘lgan yorug‘lik tushmoqda. Sirtning yuzasi $40\text{ cm}^2$ ga teng. Yorug‘likning shu sirtga beradigan bosimini ($\mu$Pa) va bosim kuchini (nN) aniqlang.
\begin{enumerate}[label=\Alph*)]
\item 1,2; 0,26.
\item 1,2; 0,13.
\item 5,6; 11,8.
\item 5,6; 22,4.
\end{enumerate}

\item Sirtga tik tushayotgan yorug‘lik 100\% qaytaruvchi sirtga 100\% yutuvchi sirtga ko‘rsatadiganiga nisbatan necha marta ko‘p bosim ko‘rsatadi?
\begin{enumerate}[label=\Alph*)]
\item ikki holda ham bir xil.
\item 1,5.
\item 2.
\item 3.
\end{enumerate}

\end{enumerate}



\begin{table}[!ht]
\centering
\begin{tabular}{|c|c|c|c|c|c|c|c|c|c|c|}
\hline
1d & 2d & 3d & 4a & 5b & 6c & 7c & 8d & 9d & 10d & 11b \\ \hline
12c & 13d & 14d & 15a & 16d & 17a & 18b & 19a & 20c & 21b & 22d \\ \hline
23d & 24c & 25b & 26d & 27b & 28d & 29c & & & & \\ \hline
\end{tabular}
\end{table}
